\documentclass[11pt]{article}
\usepackage[utf8]{inputenc}
\usepackage{amsmath}
\usepackage{geometry}
\geometry{a4paper, margin=1in}

\title{Eternal Torment: src/ Repository Overview}
\author{pale shadow}
\date{\today}

\begin{document}

\maketitle

\section{ass / ass-intel-x86}
Contains low-level assembly experimentation, focusing on AArch64 and x86-64 system calls, register operations, and build automation via GNU Autotools.

\section{palm-api}
Integration scripts for Large Language Model APIs, featuring utilities for technical documentation and Graphviz-based architectural visualizations.

\section{trithemius}
Implementation of the polyalphabetic Trithemius cipher, which utilizes a linear progressive shift logic based on character position.

\section{vigenere}
The Vigenère cipher is a method of encrypting alphabetic text by using a series of interwoven Caesar ciphers, based on the letters of a keyword. Unlike the Trithemius cipher, which uses a fixed progressive shift, the Vigenère cipher uses the characters of a secret key to determine the shift for each letter.

The mathematical formula for encryption is:
\begin{equation}
    C_i = (P_i + K_i) \pmod{26}
\end{equation}
Where $C$ is the ciphertext, $P$ is the plaintext, and $K$ is the key (repeated to match the length of the plaintext).

\section{Graphics, Math, and ML}
Includes dedicated sandboxes for:
\begin{itemize}
    \item \textbf{Graphics:} \texttt{mandelbrot}, \texttt{freeglut}, and \texttt{SpinningCube}.
    \item \textbf{Mathematics:} \texttt{calc-pi} and \texttt{projecteuler.net}.
    \item \textbf{Machine Learning:} \texttt{tensorflow} and \texttt{betty-lab}.
\end{itemize}

\end{document}
